\chapter{Using the Toolchain}
\label{ch:toolchain}

This chapter walks through building and running a program with the three tools. In this repository they are built from
the \texttt{sirc-vm} Cargo workspace; the commands below show the tool names, and \texttt{cargo run --bin \emph{tool}
	--} can be used in front of the arguments to build and run a tool in one step.

\section{A Worked Example}

Save this program as \texttt{count.sasm}. It counts \texttt{r1} up to ten and stops the simulator.

% qa:assemble
\begin{lstlisting}
    .ORG 0x0000
    .DQ @start                  ; reset vector

    .ORG 0x0200
:start
    MOVI r1, #0
:loop
    ADDI r1, #1
    CMPI r1, #10
    BRAN|!= @loop
    COPI #0x14FF                ; stop the simulator (test exit)
\end{lstlisting}

Assemble it, link it, and run it:

\begin{lstlisting}[language={}]
assembler --input-file count.sasm --output-file count.o
linker --segment-offset 0 --output-file count.bin count.o
sirc_vm --program-file count.bin --register-dump-file count.register-dump
\end{lstlisting}

The register dump written at the end shows \texttt{r1} holding \texttt{0xA}.

\section{The Assembler}

\begin{lstlisting}[language={}]
assembler --input-file <FILE> --output-file <FILE>
\end{lstlisting}

The assembler reads one source file and writes one object file; both options are required and have the short forms
\texttt{-i} and \texttt{-o}. It stops at the first syntax error and prints the offending line
(Chapter~\ref{ch:diagnostics}).

The object file is a binary file for the linker. It holds the assembled program, the labels the file declares, the
references it makes, and a copy of the source for the debug map. It is not meant to be read or edited.

\section{The Linker}

\begin{lstlisting}[language={}]
linker --segment-offset <N> --output-file <FILE> <OBJECT FILES>...
\end{lstlisting}

The linker joins its input objects in the order they are given, so the first object starts at the lowest address.
Labels declared in any object can be referred to from any other. It then patches every reference using the final
address of its label, and writes:

\begin{itemize}
	\item the program as a flat binary at the output path, and
	\item a debug map at the output path with \texttt{.dbg} added, which relates program addresses to source lines.
\end{itemize}

\texttt{--segment-offset} (short form \texttt{-s}) is added to every label's address before it is patched in. A program
that is mapped into its own memory segment, as every example in the architecture manual is, is linked with an offset of
0.

A program made of several files is linked in one step:

\begin{lstlisting}[language={}]
linker --segment-offset 0 --output-file prog.bin main.o data.o
\end{lstlisting}

\section{The Simulator}

\begin{lstlisting}[language={}]
sirc_vm --program-file <FILE> [--segment <SEGMENT>]... [--register-dump-file <FILE>]
\end{lstlisting}

The simulator loads the binary at the start of its program memory and starts the CPU through the reset vector.

\begin{description}
	\item[\texttt{--program-file}, \texttt{-p}] The linked binary to load.
	\item[\texttt{--segment}, \texttt{-s}] Maps additional memory, written \texttt{label:offset:length} or
	      \texttt{label:offset:length:file}, with the offset and length in hexadecimal. With a file, the segment is
	      initialised from it. This option may be repeated.
	\item[\texttt{--register-dump-file}, \texttt{-r}] Writes the final register state to a file when the program stops,
	      which is how the repository's examples check their results.
	\item[\texttt{-v}, \texttt{-q}] Raise or lower the amount of logging. \texttt{-vv} is common while developing.
	\item[\texttt{--debug}, \texttt{-d}] Starts a debug adapter for an editor to attach to, using the debug map written
	      by the linker.
\end{description}

For example, a program that needs 64\,K words of scratch memory away from the program binary is run with:

\begin{lstlisting}[language={}]
sirc_vm --program-file prog.bin --segment SCRATCH:00010000:FFFF
\end{lstlisting}

\section{Building with Make}

A \texttt{Makefile} keeps the three steps in sync. This is the pattern used by the examples in the repository:

\begin{lstlisting}[language={}]
prog.o: prog.sasm
	assembler --input-file prog.sasm --output-file prog.o

prog.bin: prog.o
	linker --segment-offset 0 --output-file prog.bin prog.o

run: prog.bin
	sirc_vm -vv --program-file prog.bin
\end{lstlisting}
